% GENERATED FILE. Edit canonical lessons, then run npm run generate:books.
% canonical-source-hash: fnv1a64:a7c724cc6e0e0ba7
% canonical-lessons: PA-R173-message-timed-again, PA-R173-timed-repair-again

\chapter{Return to Two Familiar Timed Writing Moves}
\label{ch:pa-timed-writing-r3}
\hlchaptermodality{173}

\begin{chapteropening}
\textbf{By the end of this chapter:} \emph{I can write a short named-reader message without a model under a practice clock, then preserve the first attempt and repair one line separately.}

Preserve the fresh timed message and write one separate Gurmukhi correction without claiming mock or readiness credit.
\end{chapteropening}

\section[(a familiar reader, a fresh short clock)]{A familiar message, a fresh first attempt}

\label{lesson:PA-R173-message-timed-again}



\begin{quote}
Close the earlier message pages. Put out a blank page and a three-minute timer. This is a fresh use of already-taught words: there is no model, romanization, word bank, dictionary, or copyable answer during the clock.
\end{quote}

\subsection*{Writing — independent timed return}
As Manan, write Aman a \textbf{30–40-word} Gurmukhi message. Greet Aman and ask how he is. Say that you like reading Punjabi and that you are well. Suggest meeting again and close in a fitting tone. Add only familiar words if you need them to reach the band; choose and order your own sentences.

Start the timer and stop after three minutes, even if the message is not finished. Record the word count and elapsed time on a separate line. Keep this first attempt unchanged for the next short lesson.

\subsection*{Your turn — observe, do not repair yet}
Without reopening an earlier page, note whether Aman is clearly the reader, the reading purpose is present, and the message is within 30–40 words. Do not rewrite the clocked page now.

\begin{tcolorbox}[breakable,colback=teal!4,colframe=teal!35!black,title={Before you move on}]
This short return is not a complete A1 mock. The later two full papers and human validation remain separate from this one timed message.
\end{tcolorbox}
\section[(one separate repair after the clock)]{Keep the clocked page; repair elsewhere}

\label{lesson:PA-R173-timed-repair-again}



\begin{quote}
Put the previous message's first attempt beside a new blank page. Keep the first attempt unchanged. Close the older examples: there is no model or copyable answer for this small repair.
\end{quote}

\subsection*{Your turn — one separate correction}
Inspect one thing at a time: named reader and purpose, sentence order, Punjabi word choice and agreement, then Gurmukhi spelling, spaces, and punctuation. Choose just one line you would improve. Write that corrected line on the new page in Gurmukhi and name the one change you made. If the first attempt has no error in that line, choose a line whose meaning could be clearer; do not invent a mistake.

\begin{tcolorbox}[breakable,colback=teal!4,colframe=teal!35!black,title={Before you move on}]
The separate line helps the next attempt; it never replaces or raises the score of the timed page. This return is not a complete A1 mock. The two full papers, calibration, and human validation remain separate gates.
\end{tcolorbox}
